\subsection{\texorpdfstring{\(C^*\) 中的可乘族}{复数乘法群中的可乘族}}

在非零复数的乘法群 \(\mathbf{C}^*\) 中，一个族 \((z_{i})_{i\in I}\) 只有当 \(\lim z_{i} = 1\) （关于 \(I\) 的有限子集的余集所成的滤子）时才可能是可乘的（第三章，§ 5，第 2 号，命题 1）；此外，由于 \(\mathbf{C}^*\) 的每一点都有可数的基本邻域系，指标 \(i\) （满足 \(z_{i}\neq 1\) 者）所成的集合当族 \((z_{i})\) 可乘时是可数的（第三章，§ 5，第 2 号，命题 1 的推论）。

命题 2 。复数的一个族 \((z_{t})\) ，其中 \(z_{t} = r_{t}(\cos \theta_{t} + i\sin \theta_{t})\) ，它是可乘的，当且仅当族 \((r_{t})\) （诸 \(z_{t}\) 的绝对值所成的）在 \(\mathbf{R}_{+}^{*}\) 中可乘，且族 \((\theta_{t})\) （诸 \(z_{t}\) 的辐角所成的）在角群 \(\mathfrak{A}\) 中可和。

鉴于群 \(\mathbf{C}^*\) 的结构（ \(\S\) 1，第 3 号，命题 1），本命题是第三章 \(\S\) 5 第 4 号命题 4 的直接推论。

若把每个角 \(\theta\) 映到它的（对任给基 \(a\) 而言的）属于区间 \([- \frac{1}{2} a, \frac{1}{2} a]\) 的那个测度，就得到 \(\mathfrak{A}\) 与 \(\mathbf{R}\) 的一个局部同构（ \(\S 2\) ，第 2 号）；由于 \(\lim \theta_{t} = 0\) （关于 \(\mathbf{I}\) 的有限子集的余集所成的滤子），我们可以把（命题 2 陈述中的）族 \(\theta_t\) 在 \(\mathfrak{A}\) 中可和这一条件，换成这样的条件：族 \((t_t)\) （诸角 \(\theta_t\) 的、属于 \([- \frac{1}{2} a, \frac{1}{2} a]\) 的测度所成之族）在 \(\mathbf{R}\) 中可和。

下述定理给出了形如 \((\mathbf{1} + u_{\iota})\) 的复数族在 \(\mathbf{C}^*\) 中可乘的另一个判别法。（它推广了第四章 § 7 第 4 号的定理 4；并参见第九章附录第 2 号命题 1：）

定理 1 。族 \((\mathbf{1} + u_{\iota})_{\iota \in \mathbf{I}}\) 在 \(\mathbf{C}^*\) 中可乘的充分必要条件是：族 \((|u_{\iota}|)\) 在 \(\mathbf{R}\) 中可和。

对 I 的每个有限子集 J，令

\[
p _ {J} = \prod_ {\iota \in J} (\mathbf{1} + a _ {\iota}), \quad s _ {J} = \sum_ {\iota \in J} a _ {\iota}, \quad \sigma_ {J} = \sum_ {\iota \in J} | a _ {\iota} |.
\]

引理 1 。对 I 的每个有限子集 J，令 \(\varphi(J) = \sup_{L \subset J} (p_L - 1)\) 。则对 J 的每个子集 L，我们有

\[
\left| p _ {\mathbf {L}} - \mathbf {1} - s _ {\mathbf {L}} \right| \leqslant \varphi (\mathrm{J}) \sigma_ {\mathbf {L}}.\tag{2}
\]

当 L 为空集时这是明显的。我们对 card (L) 作归纳。设 \(L = K \cup \{\lambda\}\) ，其中 \(\lambda \notin K\) ；则 \(p_{\mathbf{L}} = p_{\mathbf{K}}(\mathbf{1} + a_{\lambda})\) 且 \(s_{\mathbf{L}} = s_{\mathbf{K}} + a_{\lambda}\)，从而 \(p_{\mathbf{L}} - \mathbf{1} - s_{\mathbf{L}} = (p_{\mathbf{K}} - \mathbf{1} - s_{\mathbf{K}}) + (p_{\mathbf{K}} - \mathbf{1})a_{\lambda}\) ；于是由归纳假设与 \(\varphi(\mathbf{J})\) 的定义，有

\[
\left| p _ {\mathrm{L}} - \mathrm{1} - s _ {\mathrm{L}} \right| \leqslant \varphi (\mathrm{J}) \sigma_ {\mathrm{K}} + \varphi (\mathrm{J}) \left| a _ {\lambda} \right| = \varphi (\mathrm{J}) \sigma_ {\mathrm{L}},
\]

这就证明了引理。

引理 2 。若 J 是 I 的使 \(\varphi(J) < 1/4\) 的有限子集，则

\[
\left| \sigma_ {J} \right| \leqslant 4 \varphi (J) / (1 - 4 \varphi (J)).
\]

因为 \(\sigma_{\mathbf{L}} \leqslant \sigma_{\mathbf{J}}\) 对子集 \(\mathbf{L}\) （ \(\mathbf{J}\) 的每一个）成立，故由 (2) 得 \(|s_{\mathbf{L}}| \leqslant \varphi(\mathbf{J})\sigma_{\mathbf{J}} + |p_{\mathbf{L}} - \mathbf{1}| \leqslant (\mathbf{1} + \sigma_{\mathbf{J}})\varphi(\mathbf{J})\) ；但由第七章 §3 第 1 号命题 2，有 \(|\sigma_{\mathbf{J}}| \leqslant 4\sup_{\mathbf{L} \subset \mathbf{J}} |s_{\mathbf{L}}|\) ，从而 \(\sigma_{\mathbf{J}} \leqslant 4\varphi(\mathbf{J})(\mathbf{1} + \sigma_{\mathbf{J}})\) ，由此即得结论。

现在来证明定理 1 所述的条件是充分的。族 \((|u_{\iota}|)\) 在 \(\mathbf{R}\) 中可和这一假设蕴含族 \((\mathrm{1} + |u_{\iota}|)\) 在 \(\mathbf{R}_{+}^{*}\) 中可乘（第四章，§ 7，第 4 号，定理 4）；因此，对每个 \(\varepsilon > 0\) ，存在 I 的有限子集 \(J_{0}\) ，使得对每个不与 \(J_{0}\) 相交的 I 的有限子集 L，有 \(\prod_{\iota \in L} (\mathrm{1} + |u_{\iota}|) - \mathrm{1} \leqslant \varepsilon\) 。但我们可以把 \(\prod_{\iota \in L} (\mathrm{1} + u_{\iota}) - \mathrm{1}\) 写成 \(\sum_{M} \left( \prod_{\iota \in M} u_{\iota} \right)\) 的形式，其中 M 取遍 L 的所有非空子集；而由于 \(\left| \prod_{\iota \in M} u_{\iota} \right| = \prod_{\iota \in M} |u_{\iota}|\) ，我们有

\[
\left| \prod_ {\iota \in L} (1 + u _ {\iota}) - 1 \right| \leqslant \sum_ {M} \left(\prod_ {\iota \in M} | u _ {\iota} |\right) = \prod_ {\iota \in L} (1 + | u _ {\iota} |) - 1 \leqslant \varepsilon .
\]

由于 \(\mathbf{C}^*\) 是完备群，根据柯西判别法，这就证明了我们的断言。

还须证明定理 1 的条件是必要的。若 \((\mathbf{1} + u_{\iota})_{\iota\in \mathbf{I}}\) 是 \(\mathbf{C}^*\) 中的可乘族，则存在一个有限子集 \(J\) （ \(I\) 的），使得对每个有限子集 \(H\) （ \(I\) 的、不与 \(J\) 相交的），有 \(\left|\prod_{\iota\in H}(1 + u_{\iota}) - 1\right| \leqslant 1/8\) 。由引理 2 推出，\(\sum_{\iota\in H}|u_{\iota}| \leqslant 1\) 对每个有限子集 \(H\) （ \(I\) 的、不与 \(J\) 相交的）成立，因而族 \((|u_{\iota}|)\) 在 \(\mathbf{R}\) 中可和（第四章，§ 7，第 1 号，定理 1）。

上面的证明经过适当修改，也适用于某些非交换除环和代数中的（有序）无穷乘积（见习题 6 及第九章附录）。

